\documentclass[10pt,twoside]{siamart1116}
% or 
%\documentclass[10pt,twoside]{siamltex}

%Standard ELA packages
\usepackage[english]{babel}
\usepackage{graphicx,epstopdf,epsfig}
\usepackage{amsfonts,epsfig,fancyhdr,graphics, hyperref,amsmath,amssymb}

%ELA page dimensions
\setlength{\textheight}{210mm}
\setlength{\textwidth}{165mm}
\topmargin = -10mm
\renewcommand{\baselinestretch}{1.1}
\setlength{\parskip}{.1in}

% Box for end of proof outside environment
\def\cvd{~\vbox{\hrule\hbox{%
     \vrule height1.3ex\hskip0.8ex\vrule}\hrule } }

%Information for ELA papers, to be completed by the Editors upon acceptance
%Starting page
\newcommand{\SP}{1}
%Ending page
\newcommand{\EP}{10}
%Publication month
\newcommand{\PMonth}{Month}
%Handling Editor
\newcommand{\HE}{Namie of Handling Editor}
%Date of Submission
\newcommand{\DoS}{Month/Day/Year}
%Date of Acceptance
\newcommand{\DoA}{Month/Day/Year}
\newcommand{\CA}{Name of Corresponding Author}

%Please insert names of author(s) and short title for running headers
\newcommand{\Names}{Moshe Goldberg, Daniel Hershkowitz, and Daniel Szyld}
\newcommand{\Title}{Sample paper for ELA}


% Common extra environments
\newtheorem{remark}[theorem]{Remark}
\newtheorem{example}[theorem]{Example}

% Symbols for real and complex numbers
\newcommand{\reals}{\mathbb{R}}
\newcommand{\complex}{\mathbb{C}}
 
%Setting up to have Theorems indexed by section
\renewcommand{\theequation}{\thesection.\arabic{equation}}
\renewtheorem{theorem}{Theorem}[section]
%\newtheorem{remark}[theorem]{Remark}
%\newtheorem{example}[theorem]{Example}


\begin{document}

\bibliographystyle{plain}

%  Leave these commented lines here
%\input{ELAheader-template.tex}
% ELA insert correct page number
\setcounter{page}{1}

\thispagestyle{empty}

%Insert the title of the paper
 \title{\Title\thanks{Received
 by the editors on \DoS.
 Accepted for publication on \DoA. 
 Handling Editor: \HE. Corresponding Author: \CA}}

\author{
Daniel B.\ Szyld\thanks{Department of Mathematics,
Temple University, Philadelphia, Pennsylvania 19122-6094, USA
(szyld@math.temple.edu). Supported by a generous grant from ILAS.}
% Remember to put \and between any two authors
\and
Danny Hershkowitz\thanks{Mathematics Department, Technion,
Haifa 32000, Israel (hershkow@math.technion.ac.il, \newline 
goldberg@math.technion.ac.il).}
\and 
Moshe Goldberg\footnotemark[3]}


%\pagestyle{myheadings}
\markboth{\Names}{\Title}

\maketitle

\begin{abstract}
A good descriptive abstract is always very useful. The main results of
the paper are explained here of course, if possible. It should be written 
in the third person!
\end{abstract}

\begin{keywords}
Nonsingular matrices, Determinants, Spherical coordinates.
\end{keywords}
\begin{AMS}
15A15, 15F10. 
\end{AMS}

% Sample article for the Electronic Journal of Linear Algebra



%%%%%%%%%%%%%%%%%%%%%%%%%%%%%%%%%%%%%%%


%%%%%%%%%%%%%%%%%%%%%%%%%%%%%%%%%%%%%%%%%%%%%%%%%%%%%%%%%%%%%
\section{Introduction} \label{intro-sec}
Each section starts with a {\verb# \section{name} #} command.
The label is used to refer to it in other sections. Definitions, lemmas,
theorems, etc., are numbered consecutively within each section.
Formulas are numbered on the left as in the following equation:
\begin{equation} 
\label{firstfor}
Ax = b, 
\end{equation}
where $A$ is $n \times n$.

\begin{definition}
{\rm A square matrix $A$ is said to be {\em nonsingular} if its columns are 
linearly independent. }
\end{definition}

\begin{lemma} \label{main-lem}
If $A$ is $n \times n$ and nonsingular, the system {\rm (\ref{firstfor})} 
has a unique solution for any $b \in \reals^n$.
\end{lemma}

\begin{proof}
Consult any standard textbook (in linear algebra).
\end{proof}

\begin{corollary}
If $A$ is $n \times n$ and singular, the system {\rm (\ref{firstfor})} 
may have more than one solution.
\end{corollary}

\begin{theorem} \label{main-th}
Let $A$ be an $n \times n$ matrix. Then 
$\displaystyle{\lim_{k \rightarrow \infty}} A^k =O$ 
if and only if $\rho(A) <1$.
\end{theorem}

\begin{proof}
This is a well-known result; see e.g., Varga \cite[Some result]{varga-62}.
\end{proof}

%%%%%%%%%%%%%%%%%%%%%%%%%%%%%%%%%%%%%%%
\section{Another section}
A few more items worth reading before you prepare an ELA article.

\begin{remark}
 
\begin{itemize}
\item[]
\item
Remark and Definition contents are in roman font, while
Lemma and Theorem statements are in italics.
\item
Note in the bibliography how a journal article like \cite{schneider-84}
is cited.  The journal name should be either completely spelled out or,
if abbreviated, the abbreviation should be the one used by 
{\em Mathematical Reviews} ({\tt http://www.ams.org}).
Consistency should be established in each paper: either all journal 
names are abbreviated, or none.
\item
Article \cite{ela-paper} is an example of a cited ELA paper.
The official abbreviation that can be used for ELA is
{\em Electron. J. Linear Algebra}.
\item
Use the macros {\verb# \reals #} to denote the real numbers $\reals$ and 
{\verb# \complex #} to denote the complex numbers $\complex$.
\item
Square brackets {\verb# \left[#} and {\verb# \right] #} are preferred for arrays.
\item
Only formulas referenced in the text, such as 
(\ref{firstfor}) or (\ref{below}) below, should be numbered; 
others can be included in the text or be displayed without numbers. 
Similarly, one expects a reference 
in the text to every item listed in the bibliography.
\item
When the end of a proof is a displayed formula, the end of the proof
mark should be in the line of the formula. Use the {\verb# \cvd #}
macro as illustrated in Lemma \ref{lem}.
\end{itemize}
\end{remark}

\begin{lemma}
\label{lem}
The off-diagonal entries of a $2 \times 2$ symmetric matrix 
are the same.
\end{lemma}

{\em Proof.}  
The proof follows from the structure of the matrix, as follows:
\begin{equation}
\label{below}
A =  \left[ \begin{array}{lr} a  & b \\ b &  c  \end{array}
\right]. ~~~~~~ \cvd
\end{equation}

% There are various ways to include postscript figures in a
% LaTeX file. Consult the manual for details.
% Below is a particular example with an encapsulated postscript file
% Note that the epsfig package must be included in the preamble.

There are various ways to include figures.  Here are ways to include a figure 
from .eps, .pdf, and .jpg formats.

\begin{verbatim}
\begin{figure}[ht]
 \begin{center}
 \epsfig{file=imgres.png,height=40mm,width=40mm,clip=}
 \caption{A .png graphic}
 \end{center}
 \label{fig1}
 \end{figure}
\end{verbatim}


\begin{figure}[ht]
 \begin{center}
 \epsfig{file=imgres.png,height=40mm,width=40mm,clip=}
 \caption{A .png graphic}
 \end{center}
 \label{fig1}
 \end{figure}

\newpage 

\begin{verbatim}
 \begin{figure}[ht]
 \begin{center}
 \epsfig{file=logo.eps,height=25mm,width=25mm,clip=}
 \caption{An .eps graphic}
 \end{center}
 \label{fig2}
 \end{figure}
\end{verbatim}

\begin{figure}[ht]
 \begin{center}
 \epsfig{file=logo.eps,height=25mm,width=25mm,clip=}
 \caption{An .eps graphic}
 \end{center}
 \label{fig2}
 \end{figure}



\begin{verbatim}
\begin{center}
 \includegraphics[scale=.3]{unitary.pdf}
   \caption{A .pdf graphic}
\end{center}
\label{fig3}
\end{figure}
\end{verbatim}

\begin{figure}[ht]
\begin{center}
 \includegraphics[scale=.3]{unitary.pdf}
   \caption{A .pdf graphic}
\end{center}
\label{fig3}
\end{figure}

\newpage

\begin{verbatim}
\begin{figure}
\begin{center}
 \includegraphics[scale=.15]{la.jpg}
 \caption{A .jpg graphic}
\end{center}
\label{fig4}
\end{figure}
\end{verbatim}

\begin{figure}
\begin{center}
 \includegraphics[scale=.15]{la.jpg}
 \caption{A .jpg graphic}
\end{center}
\label{fig4}
\end{figure}



\bigskip
{\bf Acknowledgment.} We wish to thank everybody who has contributed
papers to the Electronic Journal of Linear Algebra.


%%%%%%%%%%%%%%%%%%%%%%%%%%%%%%%%%%%%%%%%%%%%%%%%%%%%%%%%%%%%%
\begin{thebibliography}{1}
\bibitem{schneider-84}
Hans Schneider.
\newblock  Theorems on M-splittings of a singular M-matrix which
           depend on graph structure.
\newblock  {\em Linear Algebra and its Applications}, 58:407--424, 1984.

\bibitem{varga-62}
Richard S.~Varga.
\newblock  {\em Matrix Iterative Analysis}.
\newblock  Prentice-Hall, Englewood Cliffs, New Jersey, 1962.

\bibitem{ela-paper}
S.~Friedland and H.~Schneider. 
\newblock  Spectra of expansion graphs. 
\newblock  {\em Electronic Journal of Linear Algebra}, 6:2--10, 1999.

\end{thebibliography}
\end{document}



